%% notes-tex/w037-isobutanol-scrnaseq/w037-isobutanol-scrnaseq.tex
%% [[experiments.W037-isobutanol-scrnaseq.strain-selection]]
%% https://github.com/Mjvolk3/torchcell/tree/main/notes-tex/w037-isobutanol-scrnaseq/w037-isobutanol-scrnaseq.tex
%%
%% Which of the twelve isobutanol-pathway strains on hand to grow for the first
%% single-cell RNA-seq run, and what each one is.
%%
%% Every table is generated by
%%   experiments/W037-isobutanol-scrnaseq/scripts/strain_tables.py
%% and carries a `%% SOURCE:` line naming it. Regenerate the tables before
%% rebuilding; never hand-edit anything under tables/.
%%
%% Build:  make            -> w037-isobutanol-scrnaseq.pdf       (draft view)
%%         make clean-view -> w037-isobutanol-scrnaseq-clean.pdf (share view)
%%         make check      -> formatting / provenance / style gate

\documentclass[10pt,a4paper]{article}
\usepackage[draft]{tcdoc}
%% Tables are placed with [H] so each sits under the prose that explains it.
\usepackage{float}

%% Promoter subscripts and allele superscripts. Named \psub and \gsup rather
%% than \sub and \sup because \sup is the math supremum operator and redefining
%% it would break any inline math in this document.
\newcommand{\psub}[1]{\textsubscript{#1}}
\newcommand{\gsup}[1]{\textsuperscript{#1}}

%% Our strain ids. Upright, because \gene renders italic and a strain is not a
%% gene; the paper's own YZy and SHy names are upright in this document too, so
%% the two naming schemes read alike.
\newcommand{\strain}[1]{\textup{#1}}

\title{\vspace{-1.5em}Strain Selection for the First Isobutanol
  Single-cell RNA-seq Run\\
  {\large Twelve strains on hand, six to grow}}
\author{Michael Volk}
\date{\today}

\begin{document}
\maketitle
\thispagestyle{empty}

\noindent
Twelve strains from one source paper are on hand, and all twelve are
\org{Saccharomyces cerevisiae} CEN.PK2-1C derivatives rather than S288C or
BY4741, which decides the reference genome the sequencing is mapped against.
Eight of the twelve make isobutanol. Six of those eight ferment 15\% glucose and
can therefore be grown under one protocol, and they divide into three pairs that
differ by a single plasmid, which is what makes a measured difference between
two of them attributable. Those six are the first run.
Table~\ref{tab:inventory} is everything on hand,
Table~\ref{tab:selected} is the six and
Table~\ref{tab:genotypes} their genotypes, and
Table~\ref{tab:deferred} is what is held back and why.
Section~\ref{sec:before} carries the one thing that has to be fixed first: the
only plasmid-free strain among the six needs a different growth medium than the
other five, which would otherwise confound the pair it belongs to.

\vspace{0.5em}
{\small\color{tcgray}
Tables generated by \file{experiments/W037-isobutanol-scrnaseq/scripts/strain\_tables.py}\par
Working notes: \file{notes/experiments.W037-isobutanol-scrnaseq.strain-selection.md}\par}

\vspace{1em}
\tccontents

%% notes-tex/w037-isobutanol-scrnaseq/sections/1-terms.tex
%% [[experiments.W037-isobutanol-scrnaseq.strain-selection]]
%% https://github.com/Mjvolk3/torchcell/tree/main/notes-tex/w037-isobutanol-scrnaseq/sections/1-terms.tex
%%
%% SOURCE: hand-authored. Definitions only, no data. The column names defined
%% here are produced by strain_tables.py; the numbers they describe are not
%% repeated in this section.

\section{Terms\secstatus{todo}}\label{sec:terms}

\begin{description}[leftmargin=1.2em,itemsep=2pt,topsep=2pt]

\item[add-back] A deleted gene supplied again on a plasmid. In this set the
  add-back is driven by P\psub{TDH3}, a strong constitutive promoter, so an
  add-back strain overexpresses the gene rather than restoring it to its native
  level. The distinction matters wherever an add-back is read as a stand-in for
  the wild type, and Section~\ref{sec:before} returns to it.

\item[biosensor] A two-part construct carried by all twelve strains, not only the
  isobutanol ones. The first part is the reporter, yEGFP under P\psub{LEU1},
  which Leu3p activates when bound to $\alpha$-isopropylmalate; what is sensed is
  therefore that intermediate rather than the alcohol itself. The second part is
  a leucine-insensitive $\alpha$-isopropylmalate synthase, either
  LEU4\gsup{1-410} or LEU4$\Delta$S547, which keeps producing the intermediate
  under leucine feedback. The synthase is part of the sensor rather than part of
  a production pathway.

\item[biosensor configuration] Which of three constructs a strain carries. The
  \emph{isobutanol} configuration adds a PEST degradation tag to yEGFP and pairs
  it with LEU4\gsup{1-410}, because in a \gene{LEU2}-deleted background
  $\alpha$-isopropylmalate accumulates as a by-product and raises the background
  signal. The \emph{isopentanol} configuration omits the PEST tag and pairs the
  reporter with LEU4$\Delta$S547, because \gene{LEU2} is active there and the
  intermediate is consumed rather than accumulated. A third construct is the
  reporter alone, carried by the strains that were hosts for the \gene{LEU4}
  mutant library: their synthase half arrives on the plasmid instead, so the
  working sensor is split between chromosome and plasmid. Every strain selected
  for the first run carries the isobutanol configuration, so all six express
  yEGFP and the reporter is available as a flow-cytometry check before cells are
  committed to sequencing.

\item[biosensor locus] Where the construct sits. For eleven of the twelve it is
  \gene{HIS3}, which the cassette both interrupts and restores, so a
  \gene{his3-1} parent becomes histidine prototrophic on integration.
  \strain{JC0053} is the exception and carries its biosensor at \gene{GAL80},
  marked with \texttt{natMX6}, because its \gene{HIS3} locus already holds the
  optogenetic circuit. That strain is \gene{gal80}$\Delta$ in any case, so the
  second locus costs it nothing.

\item[CEN plasmid] A centromeric plasmid, present at roughly one to two copies
  per cell, as against a 2$\mu$ plasmid at high copy. Every plasmid in this set
  is CEN and carries \gene{URA3}, so a plasmid-bearing strain is held on medium
  lacking uracil and a cell that loses the plasmid stops growing on it.

\item[$\delta$-integration] Integration into the YARCdelta5 $\delta$ sites,
  which are repeated, so the cassette lands at an unpredictable site and in an
  unpredictable number of copies. Copy number is therefore a property of the
  isolate rather than of the design, and the source paper measured it by digital
  droplet PCR, finding one to two copies in the strains it checked.

\item[high-cell-density fermentation] The production assay behind every titer
  quoted here: outgrowth in synthetic complete medium with 2\% glucose, then
  resuspension at high density in the same medium with 15\% glucose, 48\,h at
  30\,\textdegree C in a sealed 24-well plate.

\item[stated, derived] How a titer is known. A \emph{stated} titer appears in the
  source paper as a number. A \emph{derived} titer does not: the paper gives only
  a fold change against a named strain, so the absolute was obtained by dividing,
  and it is approximate and carries no uncertainty. Derived values are marked
  $\dagger$ in every table. The underlying absolutes exist only in figure bars
  and the paper's Source Data file, which is not held.

\item[titer] Product concentration in the medium at the end of a
  \emph{high-cell-density fermentation}, in mg/L. A titer is comparable to
  another only when the carbon source matches, and the carbon source is not
  constant across this set.

\end{description}

%% notes-tex/w037-isobutanol-scrnaseq/sections/2-on-hand.tex
%% [[experiments.W037-isobutanol-scrnaseq.strain-selection]]
%% https://github.com/Mjvolk3/torchcell/tree/main/notes-tex/w037-isobutanol-scrnaseq/sections/2-on-hand.tex
%%
%% SOURCE: experiments/W037-isobutanol-scrnaseq/scripts/strain_tables.py
%% Table 1 is generated by that script from the Strain records it holds. Every
%% record carries the Supplementary Table 1 row or main-text sentence it was read
%% from, plus a verbatim quote. The host genotype quoted in the prose is the
%% CEN.PK2-1C row of that same table.

\section{What is on hand\secstatus{todo}}\label{sec:onhand}

The twelve come from one paper \mirror{zhangBiosensorBranchedchainAmino2022}.
Values below were read from the primary article and its Supplementary
Information, whose sha256 are pinned in the generating script and match the
copies the literature mirror holds under that citation key, so every value here
traces to a hash-pinned artifact rather than to a live URL.

\notesec{The host is CEN.PK2-1C, not BY4741}
Supplementary Table~1 of the source paper gives the parent as
\org{S.\ cerevisiae} CEN.PK2-1C, genotype \texttt{MATa ura3-52 trp1-289
leu2-3,112 his3-1 MAL2-8c SUC2}. No strain in the paper derives from S288C or
BY4741. Two consequences follow. Reads should be mapped against a CEN.PK
assembly, because CEN.PK and S288C differ enough that an S288C-only reference
loses reads and miscalls variants. And the reference needs explicit entries for
the transgenes, since none of them are in any yeast assembly: yEGFP, the three
\gene{LEU4} alleles, \gene{Bs\_alsS}, \gene{Ec\_ilvC}\gsup{P2D1-A1},
\gene{Ll\_ilvD}, \gene{Ll\_adhA}\gsup{RE1}, \gene{EL222}, and the
\texttt{hphMX}, \texttt{kanMX}, \texttt{natMX} and \texttt{ShBle} markers.

The \gene{LEU4} alleles need care beyond being listed. A strain can be
\gene{leu4}$\Delta$ and still carry \gene{LEU4} on a plasmid, and the plasmid
allele may be full length, truncated at residue 410, or deleted for S547. Reads
from a 3\gsup{\textquotesingle}-biased protocol will not separate those alleles
unless the reference is built to separate them, and a \gene{leu4}$\Delta$ strain
carrying a plasmid copy will otherwise appear to express a gene it has lost.

%% GENERATED FILE -- do not hand-edit.
%% SOURCE: experiments/W037-isobutanol-scrnaseq/scripts/strain_tables.py
\begin{table}[H]\centering
\footnotesize
\caption[Strains on hand]{The twelve strains on hand, all of them
\org{Saccharomyces cerevisiae} CEN.PK2-1C derivatives. \emph{Titer} is the
48\,h high-cell-density value for the product named beside it; a bold value is
printed in the paper as a number, and a value marked $\dagger$ was obtained by
dividing a stated fold change, so it is approximate and carries no uncertainty.
Carbon source differs across the set and is not comparable between rows: the
three figures these numbers come from use 15\% glucose, 15\% galactose, and
2\% glucose respectively. The last column marks the six selected for the first
run.}\label{tab:inventory}
\begin{tabular}{llllll}
\toprule
id & strain & product & titer (mg/L) & carbon & first run \\
\midrule
JC0043 & YZy311 & isobutanol & $\sim$118\,$^{\dagger}$ & 15\% glucose & \checkmark \\
JC0044 & YZy313 & isobutanol & not reported & 15\% glucose & \checkmark \\
JC0045 & YZy314 & isobutanol & 1.6$\times$ JC0044 & 15\% glucose & \checkmark \\
JC0046 & YZy148 & isopentanol & not reported & 15\% glucose &  \\
JC0047 & YZy148 + LEU4 WT & isopentanol & $\sim$201\,$^{\dagger}$ & 15\% glucose &  \\
JC0048 & YZy148 + LEU4dS547 & isopentanol & \textbf{842 $\pm$ 23} & 15\% glucose &  \\
JC0049 & YZy452 & isobutanol & \textbf{310 $\pm$ 15} & 15\% galactose &  \\
JC0050 & YZy454 & isobutanol & $\sim$73\,$^{\dagger}$ & 15\% glucose & \checkmark \\
JC0051 & YZy469 & isobutanol & \textbf{227 $\pm$ 13} & 15\% glucose & \checkmark \\
JC0052 & YZy91 & isobutanol & not reported & 15\% glucose & \checkmark \\
JC0053 & YZy502 & isobutanol & not reported & 2\% glucose, dark &  \\
JC0054 & SHy134 & isopentanol & not reported & 15\% glucose &  \\
\bottomrule
\end{tabular}
\end{table}


\notesec{Four strains make isopentanol, which splits the set in two}
\strain{JC0046} through \strain{JC0048} and \strain{JC0054} produce isopentanol; the
remaining eight produce isobutanol. There is no single production ranking across
the twelve, and the highest titer on hand, 842\,$\pm$\,23\,mg/L from
\strain{JC0048}, is an isopentanol number.

\notesec{The paper's best producers are not on hand}
The strains that set the paper's production ceiling are absent from the twelve:
YZy505 at 681\,$\pm$\,29\,mg/L isobutanol, YZy148 carrying \gene{LEU4} mutant 6
at 963\,$\pm$\,32\,mg/L isopentanol, YZy470 at 443\,$\pm$\,6\,mg/L and YZy312 at
378\,$\pm$\,10\,mg/L. \strain{JC0053} is YZy502, the plasmid-free parent of
YZy505, and not YZy505 itself. Reaching the paper's ceiling means obtaining or
rebuilding strains beyond this set, which is a separate decision from the one
this run makes.

%% notes-tex/w037-isobutanol-scrnaseq/sections/3-selection.tex
%% [[experiments.W037-isobutanol-scrnaseq.strain-selection]]
%% https://github.com/Mjvolk3/torchcell/tree/main/notes-tex/w037-isobutanol-scrnaseq/sections/3-selection.tex
%%
%% SOURCE: experiments/W037-isobutanol-scrnaseq/scripts/strain_tables.py
%% Tables 2 and 3 are generated by that script. The fold changes quoted in the
%% prose are the derivation fields of the Titer records behind Table 2, and the
%% deletion lists are the deletions fields of the same records.

\section{The six to grow\secstatus{todo}}\label{sec:selection}

Six of the eight isobutanol strains ferment 15\% glucose, so one protocol covers
them. They form three pairs, and within a pair the two strains share a background
and differ only by the plasmid they carry.

%% GENERATED FILE -- do not hand-edit.
%% SOURCE: experiments/W037-isobutanol-scrnaseq/scripts/strain_tables.py
\begin{table}[H]\centering
\footnotesize
\caption[The six selected]{The six selected, as three axes of one variable
each. Within an axis the two strains share a background and differ by the
plasmid they carry, so a difference measured between them is attributable to
that plasmid. \emph{pathway} names the compartment the heterologous
isobutanol route was built in. Every strain here ferments 15\% glucose, so all
six run under one protocol. Titer conventions follow
Table~\ref{tab:inventory}.}\label{tab:selected}
\begin{tabular}{lllllll}
\toprule
axis & id & strain & deletions & pathway & plasmid & titer (mg/L) \\
\midrule
\multirow{2}{*}{1} & JC0052 & YZy91 & \textit{bat1}$\Delta$, \textit{bat2}$\Delta$, \textit{ilv6}$\Delta$ & -- & none & not reported \\
 & JC0043 & YZy311 & \textit{bat1}$\Delta$, \textit{bat2}$\Delta$, \textit{ilv6}$\Delta$ & -- & ILV6 WT & $\sim$118\,$^{\dagger}$ \\
\addlinespace
\multirow{2}{*}{2} & JC0044 & YZy313 & \textit{bat1}$\Delta$ & mitochondrial & ILV6 WT & not reported \\
 & JC0045 & YZy314 & \textit{bat1}$\Delta$ & mitochondrial & ILV6 V110E & 1.6$\times$ JC0044 \\
\addlinespace
\multirow{2}{*}{3} & JC0050 & YZy454 & \textit{ilv3}$\Delta$, \textit{tma29}$\Delta$ & cytosolic & Ll\_ilvD WT & $\sim$73\,$^{\dagger}$ \\
 & JC0051 & YZy469 & \textit{ilv3}$\Delta$, \textit{tma29}$\Delta$ & cytosolic & Ll\_ilvD I433V & \textbf{227 $\pm$ 13} \\
\bottomrule
\end{tabular}
\end{table}


\notesec{Axis 1 is a deletion against its add-back}
\strain{JC0052} is \gene{bat1}$\Delta$ \gene{bat2}$\Delta$ \gene{ilv6}$\Delta$ with
the biosensor and no pathway genes. \strain{JC0043} is that same strain carrying
\gene{ILV6} on a CEN plasmid. \gene{ILV6} encodes the regulatory subunit of
acetolactate synthase, and deleting it relieves valine inhibition, which is why
the deletion raises isobutanol in the first place. This pair is the most direct
point of contact with knockout-collection data, because the lesion is a single
gene whose isobutanol phenotype a deletion screen reports on.

\notesec{Axis 2 is one allele of one gene, in a producing background}
\strain{JC0044} and \strain{JC0045} are both YZy363, which is \gene{bat1}$\Delta$
with a five-cassette mitochondrial isobutanol pathway $\delta$-integrated and
\gene{ILV6} intact. They differ in the plasmid allele, wild type against V110E,
and the paper reports V110E as 1.6 times the wild-type titer without giving
either absolute. Because \gene{ILV6} is intact here, both strains overexpress the
gene on top of a native copy, which is a different genetic situation from Axis 1
and not a repeat of it.

\notesec{Axis 3 is a heterologous enzyme against its improved variant}
\strain{JC0050} and \strain{JC0051} are both YZy452, which is \gene{ilv3}$\Delta$
\gene{tma29}$\Delta$ with a cytosolic pathway built from
\gene{Bs\_alsS}, \gene{Ec\_ilvC}\gsup{P2D1-A1} and \gene{Ll\_ilvD}. They differ
in the plasmid allele of \gene{Ll\_ilvD}, wild type against I433V, and the
variant is 3.1 times the wild type at 227\,$\pm$\,13\,mg/L. This axis carries the
largest measured difference of the three and the highest glucose-grown titer in
the set. Its two deletions, \gene{ILV3} and \gene{TMA29}, are both genes a
deletion screen covers, so it is a second point of contact with knockout data
rather than a purely synthetic contrast.

%% GENERATED FILE -- do not hand-edit.
%% SOURCE: experiments/W037-isobutanol-scrnaseq/scripts/strain_tables.py
\begin{table}[H]\centering
\footnotesize
\caption[Genotypes]{Genotypes of the six, as Supplementary Table 1 of the
source paper gives them, with the selection markers and the growth medium each
one needs. \emph{medium} follows from the plasmid: a strain carrying a
URA3 plasmid is held on SC-ura, and JC0052, which carries none, needs uracil
instead. Section~\ref{sec:before} treats that mismatch, which is the one
thing in this list that has to be fixed before the strains are
grown.}\label{tab:genotypes}
\begin{tabular}{@{}l p{0.86\textwidth}@{}}
\toprule
id & strain and genotype \\
\midrule
\textbf{JC0052} & YZy91 \\
 & {\scriptsize \textit{his3}::HIS3-P\psub{LEU1}-yEGFP-PEST-T\psub{ADH1}-P\psub{TPI1}-LEU4\gsup{1-410}-T\psub{PGK1}, \textit{bat1}$\Delta$::hphMX, \textit{bat2}$\Delta$::lox71-kanMX-lox66, \textit{ilv6}$\Delta$::lox71-natMX-lox66} \\
 & {\scriptsize markers: hphMX, kanMX, natMX, HIS3 \quad medium: SC + uracil} \\
\addlinespace
\textbf{JC0043} & YZy311 \\
 & {\scriptsize YZy91, CEN URA3 plasmid (P\psub{TDH3}-ILV6-T\psub{ADH1})} \\
 & {\scriptsize markers: hphMX, kanMX, natMX, HIS3, URA3 \quad medium: SC-ura} \\
\addlinespace
\textbf{JC0044} & YZy313 \\
 & {\scriptsize YZy363 ($\delta$-integrated P\psub{TDH3}-ILV2, P\psub{PGK1}-ILV3, P\psub{TEF1}-CoxIVMLS-Ll\_adhA\gsup{RE1}, P\psub{TDH3}-CoxIVMLS-ARO10, P\psub{TEF1}-ILV5), CEN URA3 plasmid (P\psub{TDH3}-ILV6-T\psub{ADH1})} \\
 & {\scriptsize markers: hphMX, HIS3, ShBle, URA3 \quad medium: SC-ura} \\
\addlinespace
\textbf{JC0045} & YZy314 \\
 & {\scriptsize YZy363, CEN URA3 plasmid (P\psub{TDH3}-ILV6\gsup{V110E}-T\psub{ADH1})} \\
 & {\scriptsize markers: hphMX, HIS3, ShBle, URA3 \quad medium: SC-ura} \\
\addlinespace
\textbf{JC0050} & YZy454 \\
 & {\scriptsize YZy452, CEN URA3 plasmid (P\psub{TDH3}-Ll\_ilvD-T\psub{ADH1})} \\
 & {\scriptsize markers: hphMX, kanMX, HIS3, ShBle, URA3 \quad medium: SC-ura} \\
\addlinespace
\textbf{JC0051} & YZy469 \\
 & {\scriptsize YZy452, CEN URA3 plasmid (P\psub{TDH3}-Ll\_ilvD\gsup{I433V}-T\psub{ADH1})} \\
 & {\scriptsize markers: hphMX, kanMX, HIS3, ShBle, URA3 \quad medium: SC-ura} \\
\bottomrule
\end{tabular}
\end{table}


\notesec{What the three axes buy together}
The six span three pathway states, native flux in Axis 1, mitochondrial in
Axis 2 and cytosolic in Axis 3, three genetic backgrounds, and roughly a
threefold measured range in titer, on one carbon source and one assay format.
Axis 1 carries no engineered route at all: both of its strains make isobutanol
through endogenous valine biosynthesis and the Ehrlich pathway, with
\gene{bat1}$\Delta$ \gene{bat2}$\Delta$ keeping
$\alpha$-ketoisovalerate from being transaminated away. The plasmid that
separates its two strains carries \gene{ILV6}, a regulatory subunit rather than
a pathway enzyme, so the pair varies feedback control and not pathway capacity. Each pair isolates one variable,
so a transcriptional difference within a pair has one candidate cause. Adding
\strain{JC0044} is what makes Axis 2 a pair rather than a single high point, and
without it \strain{JC0045} has no same-background comparator anywhere in the set.

%% notes-tex/w037-isobutanol-scrnaseq/sections/4-held-back.tex
%% [[experiments.W037-isobutanol-scrnaseq.strain-selection]]
%% https://github.com/Mjvolk3/torchcell/tree/main/notes-tex/w037-isobutanol-scrnaseq/sections/4-held-back.tex
%%
%% SOURCE: experiments/W037-isobutanol-scrnaseq/scripts/strain_tables.py
%% Table 4 is generated by that script from the defer_reason field of each
%% unselected Strain record. The titers quoted in the prose are the Titer records
%% behind that same table.

\section{Held back, and why\secstatus{todo}}\label{sec:held}

%% GENERATED FILE -- do not hand-edit.
%% SOURCE: experiments/W037-isobutanol-scrnaseq/scripts/strain_tables.py
\begin{table}[H]\centering
\footnotesize
\caption[Held back]{The six held back from the first run, with the reason in
each case. Four are held back because they make isopentanol rather than
isobutanol. The other two make isobutanol but cannot share a run with the
selected six: one produces only on galactose, and one needs a light schedule.
The key below the table says how each titer is known.}\label{tab:deferred}
\begin{tabular}{llllp{0.42\textwidth}}
\toprule
id & strain & product & titer (mg/L) & reason \\
\midrule
JC0049 & YZy452 & isobutanol & \textbf{310 $\pm$ 15} & {\scriptsize galactose only: its chromosomal \textit{Ll\_ilvD} sits behind P\psub{GAL10}, so carbon source would confound every comparison in a shared run} \\
\addlinespace
JC0053 & YZy502 & isobutanol & not reported & {\scriptsize triple \textit{pdc}$\Delta$ with \gene{PDC1} restored only under a blue-light promoter, so growth and production are separate light phases and a light-shift response would sit on top of the production response} \\
\addlinespace
JC0046 & YZy148 & isopentanol & not reported & {\scriptsize makes isopentanol, not isobutanol} \\
\addlinespace
JC0047 & YZy148 + LEU4 WT & isopentanol & $\sim$201\,$^{\dagger}$ & {\scriptsize makes isopentanol, not isobutanol} \\
\addlinespace
JC0048 & YZy148 + LEU4dS547 & isopentanol & \textbf{842 $\pm$ 23} & {\scriptsize makes isopentanol, not isobutanol} \\
\addlinespace
JC0054 & SHy134 & isopentanol & not reported & {\scriptsize makes isopentanol, not isobutanol} \\
\bottomrule
\multicolumn{5}{@{}l@{}}{\scriptsize \textbf{bold} stated in the paper as a number;\quad $\dagger$ derived by dividing a stated fold change, so approximate and without uncertainty} \\
\end{tabular}
\end{table}


\notesec{JC0049 produces only on galactose}
YZy452 carries its chromosomal \gene{Ll\_ilvD} behind P\psub{GAL10}, so the
enzyme is repressed on glucose and the strain makes essentially no isobutanol
there. Its 310\,$\pm$\,15\,mg/L is the highest stated isobutanol titer on hand,
and it is a galactose number, so putting it in a glucose run would make carbon
source a second difference in every comparison involving it. Its genetics are
already represented: it is the plasmid-free parent of \strain{JC0050} and
\strain{JC0051}, which supply the same background on glucose.

\notesec{JC0053 needs a light schedule}
YZy502 is \gene{pdc1}$\Delta$ \gene{pdc5}$\Delta$ \gene{pdc6}$\Delta$ with
\gene{PDC1} restored under P\psub{C120}, which is induced by blue light, and the
first enzyme of its cytosolic pathway under a circuit that is repressed by blue
light and active in darkness. Growth and production therefore occupy different
light conditions. A single-cell measurement taken across that transition would
carry a light-shift response layered on the production response, and the paper
reports no titer for YZy502 itself, only for its plasmid-bearing derivative. The
strain is worth its own experiment with light as the variable, which is a
different design from this one.

\notesec{Four strains make the wrong product}
\strain{JC0046}, \strain{JC0047}, \strain{JC0048} and \strain{JC0054} make isopentanol.
\strain{JC0046} through \strain{JC0048} are otherwise the cleanest series in the
whole set, being one parent carrying three plasmid states across a roughly
fourfold titer range, and \strain{JC0048} holds the highest titer on hand at
842\,$\pm$\,23\,mg/L. They become the obvious first choice if isopentanol ever
takes priority.

%% notes-tex/w037-isobutanol-scrnaseq/sections/5-before-you-grow.tex
%% [[experiments.W037-isobutanol-scrnaseq.strain-selection]]
%% https://github.com/Mjvolk3/torchcell/tree/main/notes-tex/w037-isobutanol-scrnaseq/sections/5-before-you-grow.tex
%%
%% SOURCE: hand-authored, with three exceptions, each traced in place:
%%   - the empty-vector precedent is Fig. 3b and the YZy453 row of
%%     Supplementary Table 1 of the source paper;
%%   - pYZ125 is the pYZ125 row of Supplementary Table 2;
%%   - the one-to-two-copy ddPCR result is Supplementary Table 3.
%% No table in this section. The medium column it refers to is computed in
%% strain_tables.py from whether a strain carries a plasmid.

\section{Before the strains are grown\secstatus{todo}}\label{sec:before}

Three items, in the order they would block work.

\notesec{JC0052 needs an empty vector, and the source paper shows why}
Five of the six selected strains carry a \gene{URA3} CEN plasmid and are held on
SC-ura. \strain{JC0052} carries none, and the CEN.PK2-1C parent is
\gene{ura3-52}, so \strain{JC0052} needs uracil in its medium. Growing all six in
medium containing uracil would remove selection from the other five, and a CEN
plasmid is lost without selection, so by the end of a 48\,h fermentation the
population would carry plasmid-free cells. A single-cell measurement resolves
that as a distinct subpopulation, which is a confound in a run designed to
compare plasmid states.

The source paper met the same problem and solved it by transforming the
plasmid-free strain with an empty vector: Fig.~3b reports YZy91 both with and
without one, and YZy453 is YZy452 carrying empty \texttt{pYZ125}. Doing the same
here means one transformation of \strain{JC0052} with \texttt{pYZ125}, a CEN
\gene{URA3} plasmid carrying P\psub{TDH3} and a multiple cloning site and no
cargo, after which all six run on SC-ura. That strain is not among the twelve, so
the transformation has to happen before the run. Skipping it leaves Axis 1
carrying a medium difference on top of its genetic one, and the two cannot be
separated afterwards.

\notesec{Axis 1 compares a deletion to an overexpression, not to a wild type}
\strain{JC0043} supplies \gene{ILV6} from P\psub{TDH3}, so the pair
\strain{JC0052} against \strain{JC0043} is \gene{ilv6}$\Delta$ against
\gene{ILV6} overexpressed. Neither strain has \gene{ILV6} at its native level. A
comparison against knockout-collection data, where the reference is a wild type
carrying one chromosomal copy under its own promoter, has no matching strain in
this set: the strains with \gene{ILV6} intact, \strain{JC0044} and \strain{JC0045},
sit in a different background and also carry a plasmid copy. Whether that gap
matters is worth settling before the run, because closing it means building an
isogenic \gene{ILV6} wild type and that is a construction rather than a
selection.

\notesec{Transgene copy number is per-isolate on four of the six}
\strain{JC0044}, \strain{JC0045}, \strain{JC0050} and \strain{JC0051} carry
$\delta$-integrated cassettes, so their copy number was fixed when the isolate
was picked and is not specified by the construct. The source paper measured one
to two copies by digital droplet PCR in the strains it checked, which were the
sorted mitochondrial-pathway isolates rather than these four. If expression level
of a pathway gene enters the analysis, these four want their own ddPCR first,
since a twofold difference in copy number is indistinguishable from a twofold
difference in regulation once the reads are counted.


\end{document}
